% generated by roam.py handbook export -- не редактировать, править заметки
% 20261014142411-оценивающая_последовательность.org (body f86e5d3e5ab8)

\begin{definition}[{Оценивающая последовательность}]\label[definition]{note:866FE52A-109C-4912-B61C-F29BBD475F6D}
  Последовательности функций $\left\{ \varphi_k \right\}_0^{\infty}$ и чисел $\left\{ \lambda_k \right\}_0^{\infty}$, $\lambda_k \ge 0$, называется \emph{оценивающей последовательностью}
  для функции $f$, если $\lambda_k \rightarrow 0$ и
  \[
  \forall x \in \R^n\ \forall k\ \varphi_k(x) \le (1-\lambda_k) f(x) + \lambda_k \varphi_0 (x)
  \]
\end{definition}

\begin{remark}
  Каждая $\varphi_k$ не больше линейной комбинации изначальной функции $\varphi_0$ с оцениваемой функцией $f$.
\end{remark}

\paragraph*{См. также}
\begin{itemize}
  \item \hyperref[note:D21EBF43-7558-45EE-8198-60958E391E03]{Оценка невязки через оценивающую последовательность}~--- зачем она нужна
  \item \hyperref[note:C67AAFD8-D294-480B-B58E-C284BFF76D50]{Лемма о конструкции оценивающей последовательности}~--- как её построить
\end{itemize}
